% from main.tex
%%%%%%%%%%  scan idx 568  |  folio 555  %%%%%%%%%%
Montrons d'abord comment 6.2.3 découle de 6.2.2. Pour cela il suffit évidemment de démontrer (6.2.5) pour les
\(s_{(i)}=(1,\ldots,s_i,\ldots,1)\) \((i=1,\ldots,r)\), donc on peut supposer \(r=1\), \(G=\operatorname{G\ell}(n_1)\). D'après 6.2.2 \(\alpha\) se factorise par \(\operatorname{G\ell}(n_1)/\operatorname{Sl}(n_1)\) i.e. il existe un diagramme commutatif
\[
\begin{tikzcd}[column sep=large,row sep=large]
\operatorname{G\ell}(n_1) \arrow[r,"\alpha"] \arrow[d,"j"'] & \Gm\\
\operatorname{G\ell}(n_1)/\operatorname{Sl}(n_1) \arrow[ur,"\beta"'] &
\end{tikzcd}
\]
où \(j\) est l'épimorphisme canonique et \(\beta\) une représentation de \(\operatorname{G\ell}(n_1)/\operatorname{Sl}(n_1)\) dans \(\Gm\). Mais l'homomorphisme
\[
 \det:\operatorname{G\ell}(n_i)\longrightarrow\Gm
\]
possède une section, de sorte que le monomorphisme
\[
 \operatorname{G\ell}(n_i)/\operatorname{S\ell}(n_i)\longrightarrow\Gm
\]
induit par det est un isomorphisme. D'autre part on sait (SGA 3 VIII 1.3) que toute représentation de \(\Gm\) dans lui-même est de la forme
\[
 s\rightsquigarrow s^m,
\]
où \(m\) est une section de \(\mathbb Z\), d'où l'assertion.

Le corollaire 6.2.6 est une simple reformulation de 6.2.3. Notons que 6.2.3 découle aussi immédiatement de SGA 3 XXII 6.2.1, mais nous voulons lui donner ici une démonstration élémentaire via 6.2.2).

